\documentclass{article}
\usepackage[T1]{fontenc}
\usepackage[utf8]{inputenc}
\usepackage{graphicx}
\usepackage{amsmath}
\graphicspath{{../figures/}{figures/}}
\title{Growth of \textit{E. coli} stringent-response mutants}
\author{Lab}
\begin{document}
\maketitle
\section{Introduction}
The stringent response shapes bacterial growth.
\section{Results}

To quantify how the stringent-response mutations affect growth, we fitted the
integrated logistic model $y(t) = K\,[1 + e^{-r(t-t_{i})}]^{-1}$ to the
OD$_{600}$ time series of each of the three biological replicates of every
strain, treating the maximum specific growth rate $r$ and the carrying capacity
$K$ as the parameters of interest. Sigmoidal models are the standard
descriptors of bacterial growth curves, and expressing them through
biologically interpretable parameters such as $r$ and $K$ makes the estimates
directly comparable between strains and experiments
\cite{monod1949growth,zwietering1990modeling,baranyi1994dynamic}. The
three-parameter logistic described all nine curves essentially exactly
($R^{2} = 0.994$--$0.998$; root-mean-square error $0.019$--$0.037$ OD$_{600}$
units). Adding a free lower asymptote did not improve any fit (AIC was higher
in $9/9$ replicates), so no separate lag parameter was retained.

Growth rate differed strongly among the three strains
(Fig.~\ref{fig:growth}a). Wild type grew fastest, with
$r = 0.663$~h$^{-1}$ (mean; 95\% CI $0.648$--$0.678$; $n = 3$; doubling time
$1.05$~h), followed by the \textit{relA} mutant with
$r = 0.466$~h$^{-1}$ (95\% CI $0.323$--$0.610$; doubling time $1.49$~h) and
the \textit{rpoS} mutant with $r = 0.351$~h$^{-1}$ (95\% CI
$0.308$--$0.394$; doubling time $1.97$~h). Carrying capacity followed the same
rank order ($K = 1.245$, $1.181$ and $1.033$ OD$_{600}$ for wild type,
\textit{relA} and \textit{rpoS}, respectively), although the differences in
$K$ were small relative to those in $r$.

The strain effect on $r$ was significant in a one-way ANOVA
($F_{2,6} = 60.97$, $P = 1.0\times10^{-4}$, $\eta^{2} = 0.953$); the
assumptions of normality and equal variances were not violated
(Shapiro--Wilk $P = 0.084$; Levene $P = 0.238$). All three pairwise contrasts
were significant in Tukey HSD post-hoc tests with family-wise error control
(Fig.~\ref{fig:growth}b): wild type versus \textit{relA},
$+0.196$~h$^{-1}$ ($P = 0.0011$); wild type versus \textit{rpoS},
$+0.311$~h$^{-1}$ ($P = 8.6\times10^{-5}$); and \textit{relA} versus
\textit{rpoS}, $+0.115$~h$^{-1}$ ($P = 0.016$). A non-parametric
Kruskal--Wallis test supported the overall strain effect
($H_{2} = 7.20$, $P = 0.027$), but with three replicates per strain the
pairwise non-parametric contrasts have insufficient resolution to reach
significance, so the pairwise inferences rest on the ANOVA/Tukey analysis,
whose assumptions this data set satisfies.

\begin{figure}[htbp]
  \centering
  \includegraphics[width=\textwidth]{growth.pdf}
  \caption{Growth of \textit{E.\ coli} wild type and stringent-response
  mutants. \textbf{a}, OD$_{600}$ growth curves (points, mean $\pm$ s.d. of
  three biological replicates) with the integrated logistic model
  $y(t) = K\,[1 + e^{-r(t-t_{i})}]^{-1}$ fitted to each replicate (lines, mean
  of the three fitted curves). \textbf{b}, Maximum specific growth rate $r$ for
  each biological replicate (points; square with error bars, mean $\pm$ s.d.,
  $n = 3$). Significance brackets show Tukey HSD post-hoc comparisons following
  a significant one-way ANOVA ($F_{2,6} = 60.97$, $P = 1.0\times10^{-4}$):
  $^{***}P < 0.001$, $^{**}P < 0.01$, $^{*}P < 0.05$. Strains are coloured as
  in \textbf{a}.}
  \label{fig:growth}
\end{figure}

\bibliographystyle{unsrt}
\bibliography{refs}
\end{document}
